\textbf{12.} Muestre que $\pi_1(P^2,y)$ es un grupo de orden $2$, donde $P^2$ es el plano proyectivo.
\begin{proof}
La función $a:S^2\to S^2:x\mapsto-x$ es continua; como es su propia inversa, es un homeomorfismo. Tomemos $p:S^2\to P^2:x\mapsto[x]=\{x,-x\}$, que es sobreyectiva y continua por la definición de espacio cociente.\\\\
\textbf{Hecho:} Sea $[x]\in P^2$ cualquiera. Se tiene $p^{-1}([x])=\{x,-x\}$.

Como $p(x)=[x]=[-x]=p(-x)$, entonces $x,-x\in p^{-1}([x])$. Si $y\in p^{-1}([x])$ entonces $\{y,-y\}=p(y)=[x]=\{x,-x\}$, luego $y=x$ o $y=-x$. Esto prueba el hecho.\\\\
Veamos que $p$ es abierta. Sea $A\subseteqab S^2$. Tenemos
\[
p(A)=\{p(x)\mid x\in A\}=\{[x]\mid x\in A\}=\bigcup_{x\in A}\{[x]\}.
\]
Por tanto
\begin{align*}
p^{-1}(p(A))&=\bigcup_{x\in A}p^{-1}([x])\\
&=\bigcup_{x\in A}\{x,-x\}\\
&=\left(\bigcup_{x\in A}\{x\}\right)\cup\left(\bigcup_{x\in A}\{-x\}\right)\\
&=A\cup\left(\bigcup_{x\in A}\{a(x)\}\right)\\
&=A\cup a(A).
\end{align*}
Como $a$ es en particular abierta, $a(A)\subseteqab S^2$ y por tanto $p^{-1}(p(A))\subseteqab S^2$; por definición de función cociente tenemos $p(A)\subseteqab P^2$, y por tanto $p$ es abierta.\\\\
Para cada $x\in S^2$ tomamos su vecindad $V_x=B_1(x)\cap S^2$. Si existiera $y\in V_x\cap V_{-x}$ entonces $y\in B_1(x)\cap B_1(-x)$, obteniendo $2=d(x,-x)\le d(x,y)+d(y,-x)<1+1=2$. Por tanto $V_x\cap V_{-x}=\emptyset$. Notemos que si $z\in V_x=B_1(x)\cap S^2$, entonces
\begin{align*}
1&>d(x,z)\\
&=\lVert x-z\rVert\\
&=\lVert-(x-z)\rVert\\
&=\lVert(-x)-(-z)\rVert\\
&=d(-x,-z),
\end{align*}
con lo cual $-z\in B_1(-x)\cap S^2=V_{-x}$. Ahora, dado $y\in p(V_x)$, se tiene $y=p(z)$ para algún $z\in V_x$, pero $-z\in V_{-x}$, y $y=p(z)=[z]=[-z]=p(-z)$, luego $y\in p(V_{-x})$. Así, $p(V_x)\subseteq p(V_{-x})$. Como esto vale para $x\in S^2$ arbitrario, en particular $p(V_{-x})\subseteq p(V_x)$ y por tanto $p(V_x)=p(V_{-x})$.\\\\
Para cada $[x]\in P^2$ tomamos $U_{[x]}=p(V_x)$. Como $p(V_x)=p(V_{-x})$ entonces $U_{[x]}$ está bien definido, de modo que $U_{[x]}=U_{[-x]}$. Notemos que, como $x\in V_x$ entonces $[x]=p(x)\in p(V_x)=U_{[x]}$. Veamos que $p^{-1}(U_{[x]})=V_x\sqcup V_{-x}$. Como $V_x\subseteq p^{-1}(p(V_x))=p^{-1}(U_{[x]})$ y $V_{-x}\subseteq p^{-1}(p(V_{-x}))=p^{-1}(U_{[-x]})=p^{-1}(U_{[x]})$ entonces $V_x\sqcup V_{-x}\subseteq p^{-1}(U_{[x]})$. Ahora, dado $y\in p^{-1}(U_{[x]})$, se tiene $p(y)\in U_{[x]}=p(V_x)$, luego $p(y)=p(z)$ para algún $z\in V_x$, es decir $\{y,-y\}=\{z,-z\}$, de donde $y=z\in V_x$, o $y=-z\in V_{-x}$. En todo caso $y\in V_x\sqcup V_{-x}$ y por tanto $p^{-1}(U_{[x]})\subseteq V_x\sqcup V_{-x}$. De esta forma $p^{-1}(U_{[x]})=V_x\sqcup V_{-x}$, y como $V_x\subseteqab S^2$ y $V_{-x}\subseteqab S^2$, entonces $p^{-1}(U_{[x]})\subseteqab S^2$, lo cual, por definición de función cociente, implica $U_{[x]}\subseteqab P^2$.\\\\
Veamos que $p|_{V_x}:V_x\to p(V_x)$ es un homeomorfismo. Como $p$ es continua, $p|_{V_x}$ es continua. Además $p|_{V_x}$ es sobreyectiva pues $p|_{V_x}(V_x)=p(V_x)$. Si $y,z\in V_x$ y $p|_{V_x}(y)=p|_{V_x}(z)$, es decir, $\{y,-y\}=p(y)=p(z)=\{z,-z\}$, entonces $y=z$ o $y=-z$, pero con esto último $-z\in V_x$ y $z\in V_x\cap V_{-x}=\emptyset$, así que se tiene $y=z$ y por tanto $p|_{V_x}$ es inyectiva. Además $p|_{V_x}$ es abierta, pues dado $A\subseteqab V_x$, como $V_x\subseteqab S^2$ entonces $A\subseteqab S^2$, y como $p$ es abierta, $p|_{V_x}(A)=p(A)\subseteqab P^2$, y como $p|_{V_x}(A)\subseteq p(V_x)$ entonces $p|_{V_x}(A)=p|_{V_x}(A)\cap p(V_x)\subseteqab p(V_x)$. De esta forma $p|_{V_x}:V_x\to p(V_x)=U_{[x]}$ es un homeomorfismo. Como esto vale para $x\in S^2$ arbitrario, en particular $p|_{V_{-x}}:V_{-x}\to U_{[-x]}=U_{[x]}$ es un homeomorfismo.\\\\
Con lo anterior, cualquier $[x]\in P^2$ tiene una vecindad $U_{[x]}$ regularmente cubierta por $p$, donde $\{V_x,V_{-x}\}$ es la partición de $p^{-1}(U_{[x]})$ en rebanadas. Concluimos que $p:S^2\to P^2$ es una función de recubrimiento. Como $S^2$ es simplemente conexo, para cualquier $y\in P^2$ la correspondencia de levantamiento
\[
\Phi:\pi_1(P^2,y)\to p^{-1}(y)
\]
es una biyección, luego $|\pi_1(P^2,y)|=|p^{-1}(y)|=2$, es decir que $\pi_1(P^2,y)$ es un grupo de orden $2$.
\end{proof}
